\chapter{Architecture}\label{chap:architecture}

\section{Vue d'ensemble}
\soft est composé de deux couches :
\begin{description}
  \item[Cœur natif (C++)] lecture des formats, émulation des puces, mixage, sorties par voix et états des puces. Il
        est construit avec le système \texttt{make} de ZXTune pour quatre architectures Android, et produit une
        bibliothèque \texttt{libzxtune.so}.
  \item[Application (Kotlin, Java)] service de lecture, interface, listes, navigation, rendu de l'oscilloscope
        (Vulkan ou OpenGL ES) et des tableaux de bord.
\end{description}

\section{Chaîne audio}
\begin{enumerate}
  \item Un \emph{Module::Holder} est créé pour le fichier ; il fournit un \emph{Renderer} par lecture.
  \item Le renderer appelle le moteur d'émulation par blocs (environ 100~ms), rééchantillonne vers la fréquence de
        sortie (48~kHz) et applique gain, fondus et limites de boucles.
  \item Le service Android écrit le résultat dans la sortie audio. Des instantanés de l'état d'émulation, pris
        régulièrement, rendent le retour arrière rapide sans tout recalculer depuis le début.
\end{enumerate}

\section{Données de visualisation}
Chaque moteur expose les mêmes trois flux, par l'interface \emph{VoicesScope} :
\begin{description}
  \item[Disposition] liste des groupes de voix (une puce par groupe) et de leurs noms, envoyée à chaque changement.
  \item[Sorties par voix] échantillons de chaque voix, au rythme de la sortie, rééchantillonnés comme le son.
  \item[États des voix] instantané d'environ 1~ms de chaque voix : fréquence, niveau, indicateurs, type de puce et
        64 octets propres à la puce (registres, enveloppes, table d'ondes\dots).
\end{description}
Pour les puces de libvgm, ces flux viennent d'une extension ajoutée aux cœurs d'émulation : une fonction de rappel
par échantillon (sorties par voix) et une fonction d'état. Les cœurs reSID et reSIDfp, ares et Ootake ont reçu des
extensions équivalentes.

\section{Rendu}
\begin{itemize}
  \item Un fil de calcul prépare chaque image : pour chaque voix, la fenêtre de temps est calée sur la période
        fondamentale de la note (déclenchement), puis réduite à deux points par colonne de pixels (minimum et
        maximum), ce qui ne perd aucun pic.
  \item Le rendu des traces est fait par Vulkan quand il est disponible, sinon par OpenGL ES, avec un tracé fin
        anticrénelé.
  \item Le tableau de bord est dessiné en deux couches : une couche fixe (cadres, légendes), redessinée seulement
        quand la disposition change, et une couche dynamique (valeurs, courbes), redessinée à chaque image.
\end{itemize}

\section{Tests}
\begin{itemize}
  \item Plus de 1~000 tests JVM (Robolectric) pour l'application : modèles, stockage des listes, calculs des vues,
        et tests de rendu qui produisent les captures de ce manuel.
  \item Outils natifs de contrôle : mesure de l'oscilloscope (\texttt{scope\_bench}), vérification des sorties
        par voix des puces VGM, des paramètres SID, de la recherche dans le morceau.
\end{itemize}
